\Needspace{6\baselineskip}
\begin{center}\hypertarget{native-254-article181}{}\pdfbookmark[2]{181.}{article-181}181.\end{center}
Thus whenever there is only one unknown $p$, and the values of the functions $ap+n$, $a'p+n'$, $a''p+n''$, etc. that serve to determine it have been found to be $=M$, $M'$, $M''$, etc., respectively, by observations of equal accuracy, the most probable value of $p$ will be
\[
=\frac{am+a'm'+a''m''+\text{etc.}}{aa+a'a'+a''a''+\text{etc.}}=A,
\]
where we write $m$, $m'$, $m''$, etc. for $M-n$, $M'-n'$, $M''-n''$, etc., respectively.

To estimate the degree of precision to be expected in this value, we shall suppose that the probability of an error $\Delta$ in the observations is expressed by $\frac{h}{\BareRoot{\pi}}e^{-hh\Delta\Delta}$. Hence the probability that the true value of $p$ is $=A+p'$ will be proportional to the function
\[
e^{-hh((ap-m)^2+(a'p-m')^2+(a''p-m''^2+\text{etc.})},
\]
when $p$ is replaced by $A+p'$. The exponent of this function can be reduced to the form $-hh(aa+a'a'+a''a''+\text{etc.})(pp-2pA+B)$, where $B$ is independent of $p$. The function itself will therefore be proportional to
\[
e^{-hh(aa+a'a'+a''a''+\text{etc.})p'p'}.
\]
It is thus clear that the value $A$ should be assigned the same degree of precision as if it had been found by a direct observation whose precision was to that of the original observations as $h\BareRoot{(aa+a'a'+a''a''+\text{etc.})}$ to $h$, or as $\BareRoot{(aa+a'a'+a''a''+\text{etc.})}$ to $1$.
